\section{Introduction}
\label{sec:intro}


Agent \emph{harness}—the code that wraps an LLM into a working agent—has become a critical factor in agent performance: two agents built on the same LLM but with different harnesses can differ by as much as 6 $\times$ on the same benchmark~\cite{tian2026swe}. This gap has motivated \emph{automated harness optimization}: the outer-loop systems  where an LLM-driven \emph{proposer} iteratively proposes candidate agents, prompts, or memory components, and an \emph{evaluator} evaluates them on a benchmark, and uses the results to guide the next proposal~\cite{romeraparedes2024funsearch,novikov2025alphaevolve,sharma2025openevolve,agrawal2025gepa,lee2026metaharness}. 


Most prior approaches to automated harness optimization adapt techniques from \emph{text optimization}~\cite{hu2024adas,sharma2025openevolve,zhang2025ace}, in which the proposer iteratively refines a text or code artifact using feedback from prior runs. However, the context limits force this feedback to be heavily compressed into scalar scores~\cite{novikov2025alphaevolve,agrawal2025gepa}, short summaries~\cite{lee2025feedback}, or only the most recent candidate~\cite{madaan2023self,yang2023large}, which inevitably remove information needed to trace downstream failure to earlier harness decision~\cite{lee2026metaharness}.  To address this limitation, Meta-Harness~\cite{lee2026metaharness} replaces the proposer with a coding agent — an LLM equipped with developer tools — and exposes the full revision history of prior candidates as a navigable filesystem. The coding agent, named meta-agent, reads raw code and execution traces directly, via standard operations such as \textit{grep} and \textit{cat}, and decides for itself which prior attempts to inspect before proposing the next edit. This shift from compressed feedback to argentic, on-demand inspection has produced clear gains over various baselines~\cite{lee2026metaharness}.

Yet full access to the revision history is not the same as effective use of it. A natural question is whether the meta-agent can reliably locate what matters for the next edit: as the optimization horizon grows, the volume of available context may dilutes the meta-agent's attention rather than informing it. The agent must repeatedly re-discover, within an ever-expanding pool of prior code and execution traces, which subset is actually relevant to the current failure. The question is therefore not "what should the meta-agent be \emph{allowed to see}?" but "what should the meta-agent \emph{attend to}?" In this paper, we reframe automated harness optimization as a \emph{context curation} problem and ask: 
\begin{quote}
What is the optimal context to expose to the meta-agent at each round?
\end{quote}

To answer this question, we introduce \textbf{\name}, a framework that treats context curation as the central optimization variable in automated harness optimization. \name sits between the optimization loop and the meta-agent, deciding at each round which past iterations and which files the meta-agent should attend to. Within this framework, we develop two complementary curation strategies that target the iteration axis and the file axis.

\textbf{Iteration-level: an adaptive budget.} We define three budget levels. The \emph{low} budget exposes only the best harness seen so far; the \emph{medium} budget adds the worst, providing a contrast that highlights which design choices matter; the \emph{high} budget exposes the full history and lets the meta-agent identify useful files itself. We start every iteration at the low budget and escalate to medium and then high only if the meta-agent fails to produce an improving edit, returning to the low budget once it succeeds. This yields cheap context on average and rich context exactly when the agent needs it.

\textbf{File-level: a bandit over files.} We treat each file in the harness as an arm of a multi-armed bandit; at every round, the bandit selects a subset of files to expose, with reward drawn from the validation improvement observed when those files were included. Over rounds, the bandit concentrates the meta-agent's attention on files that historically drive productive edits and phases out files that consistently do not. Unlike a fixed top-$k$ retriever, the bandit adapts as the harness evolves and as the locus of failure shifts across files

We evaluate \name on three benchmarks spanning agentic software engineering and long-horizon memory: SWE-bench~\cite{jimenez2024swebench}, LongMemEval~\cite{wu2024longmemeval}, and LoCoMo~\cite{maharana2024locomo}. Across all three, \name improves over the uncurated Meta-Harness baseline by [XX]\% relative on average while reducing meta-agent input length by [YY]\% (Section~\ref{...}). The two curation axes are complementary rather than redundant: iteration-only curation yields [+a]\%, file-only curation yields [+b]\%, and combining them adds a further [+c]\% (Section~\ref{...}).


% LLM-based agents are now being applied to a widening set of tasks including software engineering~\cite{jimenez2024swebench,yang2024sweagent}, memory-augmented assistants~\cite{wu2024longmemeval}, completing operating-system task~\cite{} and so on.
% Behind each of these systems is a \emph{harness}: the prompts, tool definitions, retrieval rules, and control flow that turn a base LLM into a goal-directed agent~\cite{lee2026metaharness}. Recently, hand-engineering this harness has emerged as a primary bottleneck in agent development, motivating a recent wave of work on \textbf{automated harness optimization}, in which a \emph{meta-agent} iteratively edits the harness over many rounds, guided by validation-set feedback and a record of prior attempts~\cite{hu2024adas,sharma2025openevolve,zhang2025ace,lee2026metaharness}

% This optimization loop is, however, more fragile than its mechanics suggest. By round $N$, the optimization record consists of every prior harness, the diffs between them, validation traces from each round, and the agent's full  source tree-often often exceeding context window of a LLM. To alleviate the issue, existing methods compress iteration history into summaries~\cite{sharma2025openevolve} or curated playbooks~\cite{zhang2025ace}, feeding the meta-agent a lossy but compact view. 


% In modern pipelines this record is rarely 
% flattened into a single prompt; the meta-agent is itself an agent—often realized with Claude Code or a similar CLI-equipped backbone~\citep{...}—that decides for itself which files and which past iterations to read. The choice of what to read is therefore made by the meta-agent every round, but it is made \emph{unsupervised}: existing automated harness optimization pipelines neither curate this input nor train the meta-agent to curate it well. What the meta-agent ends up attending to is left to chance, and is rarely studied as a design variable in its own right. Meta-Harness~\cite{lee2026metaharness} takes the opposite stance,
% deploying a coding agent with full file system access to all prior candidates, scores, and traces, on the premise that a capable agent's native browsing and reasoning already constitute an effective search gradient. 




% Large language model (LLM) agents are increasingly deployed as autonomous
% software engineers~\cite{jimenez2024swebench,yang2024sweagent}, long-context
% assistants~\cite{wu2024longmemeval}, and conversational
% companions~\cite{maharana2024locomo}. Their effectiveness is largely
% determined by \emph{harness engineering}: the design of the surrounding
% code that orchestrates memory, retrieval, prompt routing, and tool
% use~\cite{lee2026metaharness}. Automating this design through
% propose--evaluate loops, where a second LLM (the \emph{proposer})
% iteratively improves the harness, has emerged as a promising
% paradigm~\cite{hu2024adas,khattab2023dspy}.

% Existing approaches span a spectrum of information exposure.
% ACE~\cite{zhang2025ace} and OpenEvolve~\cite{sharma2025openevolve}
% compress iteration history into summaries or curated playbooks, feeding
% the proposer a lossy but compact view.
% Meta-Harness~\cite{lee2026metaharness} takes the opposite stance,
% deploying a coding agent with full filesystem access to all prior
% candidates, scores, and traces, on the premise that a capable agent's
% native browsing and reasoning already constitute an effective search
% gradient. This intuition is sound: a coding agent's chain-of-thought
% \emph{is} the optimization gradient, and its ability to autonomously
% navigate large workspaces makes it a natural proposer. Yet even powerful
% proposers benefit from active management of which files they attend to.
% As harnesses grow in complexity, unrestricted file access can dilute the
% proposer's \textbf{attention} and diffuse its \textbf{optimization
% direction}, analogous to how CODEFILTER~\cite{li2025codefilter} shows
% that filtering retrieved context to only positive-impact chunks improves
% code completion over providing all retrieved files. Explicit
% mechanism-level guidance further enables \textbf{directed exploration},
% steering the proposer toward structurally distinct alternatives rather
% than repeatedly revisiting familiar axes such as prompt wording.

% This paper introduces \textbf{Optimization Harness}, a modular framework
% that wraps any propose--evaluate loop with three capabilities:
% (i)~Docker file-scope sandboxing that restricts the proposer to a
% curated workspace, (ii)~mechanism-level optimization-focus prompts that
% steer exploration along named axes, and (iii)~labeled best/worst
% reference iterations that provide imitation targets and anti-patterns.
% Two adaptive context-selection policies build on this framework: a
% \emph{progressive} policy that escalates a budget state machine
% (\texttt{low}$\,\to\,$\texttt{medium}$\,\to\,$\texttt{high}) only after
% a Pareto-frontier stall, and a \emph{bandit} policy that scores each
% workspace file with a per-file UCB scorer and concentrates high-utility
% files in a curated ``hot workspace'' each iteration.

% Across three benchmarks (LoCoMo, LongMemEval, SWE-bench Verified) and
% two proposer families (Claude Code with Kimi K2.6, and Codex GPT-5.4),
% the progressive policy outperforms the fixed-high default in five of six
% cells. The bandit policy achieves the single strongest result, lifting
% SWE-bench Verified from 0.44 to 0.64 (+20\,pp) on the full 500-task
% split. Across ten adaptive runs, 35 of 48 score-improvement events occur
% at \texttt{low} or \texttt{medium} budget, confirming that
% \emph{managing what the optimizer sees} is a complementary and effective
% lever to improving the optimizer itself. Our contributions are:
% \begin{enumerate}
%     \item We formalize \emph{proposer context management} in harness
%     optimization and show that curating file access improves attention
%     focus, optimization direction, and exploration diversity.
%     \item We introduce Optimization Harness, a modular framework with
%     Docker sandboxing, mechanism-focus prompts, and reference labeling
%     that enables pluggable context-selection policies.
%     \item We propose two adaptive policies (progressive and bandit) and
%     evaluate them on three benchmarks with two proposer families,
%     demonstrating consistent improvements over the fixed-high baseline.
%     \item We provide an empirical analysis of where breakthroughs occur
%     across budget tiers, showing that curating attention is more
%     efficient per opportunity than maximizing context.
% \end{enumerate}